% ============================================================================
% MẪU CẤU TRÚC NGÂN HÀNG CÂU HỎI XÁC SUẤT — LỚP 9, CHƯƠNG 8
% Dùng chung với các khai báo trong main.tex và gói ex_test.
% File này là mẫu để sao chép; hiện chưa được \input vào main.tex.
% ============================================================================

% ----------------------------------------------------------------------------
% I. CÁCH SỬ DỤNG
% ----------------------------------------------------------------------------
% 1. Sao chép file này và đổi tên theo dạng bài, nên đặt trong thư mục Data.
% 2. Thay X trong "Dạng X", tên dạng bài và tất cả ID bằng số dạng từ 1 đến 8.
% 3. Chèn từng câu vào đúng mức độ và đúng loại câu hỏi.
% 4. Trước khi thêm, tìm trong cả 8 file Data để tránh câu trùng hoặc chỉ đổi tên.
% 5. Mỗi câu phải có:
%       - môi trường \begin{ex} ... \end{ex};
%       - ID ngay sau \begin{ex};
%       - đúng một khối \loigiai{...};
%       - đáp án đúng được đánh dấu bằng \True.
% 6. Sau khi hoàn thành, thêm vào main.tex:
%       \input{Data/TEN_FILE_KHONG_CO_DUOI_TEX}

% ----------------------------------------------------------------------------
% II. QUY ƯỚC ID
% ----------------------------------------------------------------------------
% Cấu trúc: %[9C8Mi-XL]
%
%   9  : lớp 9
%   C8 : chương 8
%   M  : mức độ, gồm N, H, V, C
%        N = Nhận biết; H = Thông hiểu; V = Vận dụng; C = Vận dụng cao
%   i  : nhóm bài 25–26 đang dùng trong ngân hàng hiện tại
%   X  : số dạng bài, từ 1 đến 8
%   L  : loại câu hỏi, gồm T, D, S
%        T = Trắc nghiệm; D = Đúng sai; S = Trả lời ngắn
%
% Ví dụ:
%   %[9C8Ni-1T] : Nhận biết, Dạng 1, Trắc nghiệm
%   %[9C8Hi-4D] : Thông hiểu, Dạng 4, Đúng sai
%   %[9C8Vi-6S] : Vận dụng, Dạng 6, Trả lời ngắn
%   %[9C8Ci-8T] : Vận dụng cao, Dạng 8, Trắc nghiệm
%
% ID hiện tại là mã phân loại, vì vậy nhiều câu trong cùng một nhóm có thể dùng
% cùng một ID. Không tự ý đổi cấu trúc ID nếu chưa sửa quy ước của toàn bộ dự án.

% ----------------------------------------------------------------------------
% III. DANH MỤC 8 DẠNG BÀI
% ----------------------------------------------------------------------------
% Dạng 1: Viết tập hợp không gian mẫu của một phép thử.
% Dạng 2: Xác định số phần tử của không gian mẫu của một phép thử.
% Dạng 3: Tính xác suất trong trò chơi tung đồng xu.
% Dạng 4: Tính xác suất trong trò chơi gieo xúc xắc.
% Dạng 5: Tính xác suất trong trò chơi vòng quay số.
% Dạng 6: Tính xác suất trong trò chơi rút thẻ, rút bóng từ hộp hoặc túi.
% Dạng 7: Tính xác suất trong phép thử viết hoặc chọn số tự nhiên.
% Dạng 8: Tính xác suất khi chọn, xếp một hoặc nhiều đối tượng trong nhóm.

% ----------------------------------------------------------------------------
% IV. GỢI Ý PHÂN LOẠI MỨC ĐỘ
% ----------------------------------------------------------------------------
% N — Nhận biết:
%   Nhắc lại định nghĩa, kí hiệu, nhận diện kết quả, liệt kê trực tiếp ít trường hợp.
% H — Thông hiểu:
%   Giải thích được cách lập không gian mẫu, phân biệt thứ tự, đếm bằng bảng đơn giản.
% V — Vận dụng:
%   Kết hợp nhiều điều kiện; xét chẵn, lẻ, nguyên tố, chia hết; xử lí có/không hoàn lại.
% C — Vận dụng cao:
%   Nhiều bước suy luận, nhiều trường hợp, dữ kiện ẩn hoặc cần phát hiện lỗi/ngoại lệ.

% ----------------------------------------------------------------------------
% V. MẪU CÂU HỎI — SAO CHÉP KHỐI CẦN DÙNG
% ----------------------------------------------------------------------------

% MẪU T — TRẮC NGHIỆM BỐN PHƯƠNG ÁN
%
% \begin{ex}%[9C8Ni-XT]
%   Nội dung câu hỏi là:
%   \choice
%   {Phương án A}
%   {\True Phương án B đúng}
%   {Phương án C}
%   {Phương án D}
%   \loigiai{
%     Trình bày lời giải, kiểm tra từng điều kiện và kết luận chọn B.
%   }
% \end{ex}

% MẪU D — ĐÚNG SAI BỐN Ý
%
% \begin{ex}%[9C8Hi-XD]
%   Nội dung chung của câu hỏi. Xét tính đúng sai của các khẳng định sau:
%   \choiceTF
%   {\True Khẳng định a đúng}
%   {Khẳng định b sai}
%   {\True Khẳng định c đúng}
%   {Khẳng định d sai}
%   \loigiai{
%     Nêu không gian mẫu hoặc bảng kết quả trước khi xét từng ý.
%     \begin{itemchoice}
%       \itemch \textbf{Đúng.} Giải thích ý a.
%       \itemch \textbf{Sai.} Giải thích ý b và nêu kết quả đúng.
%       \itemch \textbf{Đúng.} Giải thích ý c.
%       \itemch \textbf{Sai.} Giải thích ý d và nêu kết quả đúng.
%     \end{itemchoice}
%   }
% \end{ex}

% MẪU S — TRẢ LỜI NGẮN
%
% \begin{ex}%[9C8Vi-XS]
%   Nội dung câu hỏi.
%   \shortans{Đáp án ngắn}
%   \loigiai{
%     Trình bày lời giải và kết luận đúng bằng đáp án đã khai báo.
%   }
% \end{ex}
%
% Có thể điều chỉnh độ rộng ô đáp án, ví dụ:
%   \shortans[6]{$\dfrac{3}{10}$}
%   \shortans[10]{$\{SS;SN;NS;NN\}$}

% ----------------------------------------------------------------------------
% VI. KIỂM TRA TRƯỚC KHI LƯU CÂU HỎI
% ----------------------------------------------------------------------------
% - Câu hỏi xác định rõ kết quả được ghi nhận là cặp, bộ, tổng, tích hay màu sắc.
% - Nêu rõ có xét thứ tự hay không; có hoàn lại đối tượng đã rút hay không.
% - Chỉ dùng công thức xác suất cổ điển khi các kết quả đang đếm đồng khả năng.
% - Không nhầm số cặp kết quả với số giá trị tổng hoặc tích khác nhau.
% - Với câu đúng sai, \choiceTF có đúng 4 ý và \loigiai có đúng 4 \itemch.
% - Với câu trắc nghiệm, chỉ có một phương án được gắn \True.
% - Với câu trả lời ngắn, đáp án trong \shortans phải khớp kết luận của lời giải.
% - Đối chiếu phép tính, đáp án nguồn và các trường hợp biên như số 0, số 1.
% - Nếu có hình, ưu tiên dựng bằng TikZ và giữ đúng dữ kiện của hình gốc.

% ============================================================================
% DẠNG X: THAY BẰNG TÊN DẠNG BÀI
% Lớp 9 | Chương 8 | Bài 25–26 (i) | Mức độ: N, H, V, C | Loại: T, D, S
% Đối chiếu nội dung trong cả 8 file trước khi chèn; không chèn câu trùng.
% Mỗi câu dùng môi trường ex, gắn ID ngay sau \begin{ex} và có \loigiai{...}.
% ============================================================================
\section*{Dạng X. Thay bằng tên dạng bài}
\setcounter{ex}{0}

% ============================================================================
% MỨC ĐỘ N — NHẬN BIẾT
% ============================================================================
\subsection*{I. Nhận biết}

% ----------------------------------------------------------------------------
% TRẮC NGHIỆM (T) | ID: %[9C8Ni-XT]
% Chèn các câu dùng \choice tại đây.
% ----------------------------------------------------------------------------
\subsubsection*{1. Trắc nghiệm}


% ----------------------------------------------------------------------------
% ĐÚNG SAI (D) | ID: %[9C8Ni-XD]
% Chèn các câu dùng \choiceTF tại đây.
% ----------------------------------------------------------------------------
\subsubsection*{2. Đúng sai}

% ----------------------------------------------------------------------------
% TRẢ LỜI NGẮN (S) | ID: %[9C8Ni-XS]
% Chèn các câu dùng \shortans tại đây.
% ----------------------------------------------------------------------------
\subsubsection*{3. Trả lời ngắn}

% ============================================================================
% MỨC ĐỘ H — THÔNG HIỂU
% ============================================================================
\subsection*{II. Thông hiểu}

% ----------------------------------------------------------------------------
% TRẮC NGHIỆM (T) | ID: %[9C8Hi-XT]
% Chèn các câu dùng \choice tại đây.
% ----------------------------------------------------------------------------
\subsubsection*{1. Trắc nghiệm}

% ----------------------------------------------------------------------------
% ĐÚNG SAI (D) | ID: %[9C8Hi-XD]
% Chèn các câu dùng \choiceTF tại đây.
% ----------------------------------------------------------------------------
\subsubsection*{2. Đúng sai}

% ----------------------------------------------------------------------------
% TRẢ LỜI NGẮN (S) | ID: %[9C8Hi-XS]
% Chèn các câu dùng \shortans tại đây.
% ----------------------------------------------------------------------------
\subsubsection*{3. Trả lời ngắn}

% ============================================================================
% MỨC ĐỘ V — VẬN DỤNG
% ============================================================================
\subsection*{III. Vận dụng}

% ----------------------------------------------------------------------------
% TRẮC NGHIỆM (T) | ID: %[9C8Vi-XT]
% Chèn các câu dùng \choice tại đây.
% ----------------------------------------------------------------------------
\subsubsection*{1. Trắc nghiệm}

% ----------------------------------------------------------------------------
% ĐÚNG SAI (D) | ID: %[9C8Vi-XD]
% Chèn các câu dùng \choiceTF tại đây.
% ----------------------------------------------------------------------------
\subsubsection*{2. Đúng sai}

% ----------------------------------------------------------------------------
% TRẢ LỜI NGẮN (S) | ID: %[9C8Vi-XS]
% Chèn các câu dùng \shortans tại đây.
% ----------------------------------------------------------------------------
\subsubsection*{3. Trả lời ngắn}

% ============================================================================
% MỨC ĐỘ C — VẬN DỤNG CAO
% ============================================================================
\subsection*{IV. Vận dụng cao}

% ----------------------------------------------------------------------------
% TRẮC NGHIỆM (T) | ID: %[9C8Ci-XT]
% Chèn các câu dùng \choice tại đây.
% ----------------------------------------------------------------------------
\subsubsection*{1. Trắc nghiệm}

% ----------------------------------------------------------------------------
% ĐÚNG SAI (D) | ID: %[9C8Ci-XD]
% Chèn các câu dùng \choiceTF tại đây.
% ----------------------------------------------------------------------------
\subsubsection*{2. Đúng sai}

% ----------------------------------------------------------------------------
% TRẢ LỜI NGẮN (S) | ID: %[9C8Ci-XS]
% Chèn các câu dùng \shortans tại đây.
% ----------------------------------------------------------------------------
\subsubsection*{3. Trả lời ngắn}